%%%%%%%%%%%%%%%%%%%%%%%%%%%%%%%%%%%%%%%%%
% Medium Length Professional CV
% LaTeX Template
% Version 2.0 (8/5/13)
%
% This template has been downloaded from:
% http://www.LaTeXTemplates.com
%
% Important note:
% This template requires the resume.cls file to be in the same directory as the
% .tex file. The resume.cls file provides the resume style used for structuring the
% document.
%
%%%%%%%%%%%%%%%%%%%%%%%%%%%%%%%%%%%%%%%%%

%----------------------------------------------------------------------------------------
%	PACKAGES AND OTHER DOCUMENT CONFIGURATIONS
%----------------------------------------------------------------------------------------
\documentclass{resume} % Use the custom resume.cls style

\usepackage{amsmath}
%\usepackage[left=0.75in,top=0.44in,right=0.75in,bottom=0.44in]{geometry} % Document margins
%\usepackage{lipsum} % for dummy text
%\usepackage{enumitem}
%\setlist{nosep}
\usepackage{bibentry}

\usepackage{graphicx}
\usepackage{comment}

\usepackage{enumitem}
\usepackage{xcolor} % color
\definecolor{cvLink}{HTML}{2D4F70}
\usepackage[
  colorlinks=true,
  urlcolor=cvLink,
  linkcolor=cvLink,
  citecolor=cvLink
]{hyperref}
\usepackage{pifont}

\usepackage{CJKutf8} % Japanese
\usepackage{float}

%http://mirrors.ibiblio.org/CTAN/fonts/fontawesome/doc/fontawesome.pdf
%https://artofproblemsolving.com/wiki/index.php/LaTeX:Symbols
%usepackage for symbols eg, cell , mail
\usepackage{fontawesome} % symbols, e.g., github
\usepackage{amssymb}
%\usepackage{fontawesome5}


\usepackage{geometry}
\geometry{a4paper,left=1cm,right=1cm,top=1.3cm,bottom=1.3cm}

%\newcommand{\tab}[1]{\hspace{.2667\textwidth}\rlap{#1}}
\newcommand{\itab}[1]{\hspace{0em}\rlap{#1}}
\newcommand\tab[1][0.5cm]{\hspace*{#1}}
\newcommand{\parsec}{\texttt{PaRSEC}}
\newcommand{\dplasma}{\texttt{DPLASMA}}
\newcommand{\awards}{\texttt{Awards}}
\newcommand{\ecp}{\texttt{ECP}}
\renewcommand{\labelitemi}{$\diamond$}
\newcommand{\hicma}{\texttt{HiCMA}}
\newcommand{\openmpi}{\texttt{Open MPI}}


\let\OrigRSection\rSection
\let\OrigEndRSection\endrSection
\let\OrigMakeUppercase\MakeUppercase

\renewenvironment{rSection}[1]{
    \renewcommand{\MakeUppercase}[1]{\textcolor{blue}{\OrigMakeUppercase{##1}}}%
    \OrigRSection{#1}%
}{%
    \OrigEndRSection%
}


%\diamond
\name{Qinglei Cao} % Your name
\address{
\faEnvelopeO~\href{mailto:qingcao@iu.edu}{qingcao@iu.edu}
\qquad \faEdge~\href{https://www.qingleicao.com}{www.qingleicao.com}}
% Your address
\vspace{+2ex}


\begin{document}
\begin{rSection}{\faBookmark ~Research INTERESTS}
  %High Performance Computing (HPC) and Artificial Intelligence (AI): task-based runtime system, linear algebra algorithm, extreme-scale climate and weather modeling, and large-scale machine learning \& deep learning
  %High Performance Computing (HPC) and Artificial Intelligence (AI): parallel and distributed computing, task-based runtime system, linear algebra algorithm, extreme-scale climate and weather modeling, and large-scale machine learning \& deep learning
  %Parallel and distributed computing, Task-based runtime systems, Linear algebra algorithms, Quantum simulation, Large-scale machine learning \& deep learning, and Extreme-scale scientific applications
  High-Performance Computing (HPC), Artificial Intelligence (AI), and Quantum Computing (QC): Task-Based Runtime Systems, Numerical Linear Algebra, Mixed-Precision Computing, Quantum Circuit Simulation, Large-Scale Machine Learning \& Deep Learning, and Extreme-Scale Scientific Computing.
\end{rSection}

%----------------------------------------------------------------------------------------
%	EDUCATION SECTION
%----------------------------------------------------------------------------------------

\begin{rSection}{\faGraduationCap~Education}

  {\bf The University of Tennessee, Knoxville (UTK)}, Computer Science \hfill {2016-2022}
  \\ {\bf PhD}, High Performance Computing %\hfill {\em Expected graduation July 2022}
  %\\ {\bf Advisor}: Dr. Jack Dongarra ({Turing Award 2021})
  \\ {\bf Advisor}: Dr. Jack Dongarra ({\color{red} Turing Award 2021})
  \\ {\bf Group Leader \& Co-Advisor}: Dr. George Bosilca
  %\\ {\bf Awards}: {\bf ACM Gordon Bell Prize Finalist}, Graduate Student Senate Travel Awards, and SIAM Student Travel Award
  \vspace{+1.5ex}
  \\{\bf Ocean University of China (OUC)}, Computer Application Technology  \hfill {2013-2016}
  \\ {\bf M.E.}, Image Processing \& Parallel Computing
  %\\ {\bf Advisors}: Dr. Yuntao Qian (Zhejiang University), Dr. Zhiqiang Wei (OUC)
  %\\ {\bf Awards}: Honor of “Outstanding Graduates” (2014); Graduate Student Scholarship (2014)
  \vspace{+1.5ex}
  \\{\bf Hunan University (HNU)}, Information and Computational Science \hfill {2005-2009}
  \\ {\bf BS}, Mathematics
  %\\ {\bf Awards}: Honor of “Excellent Student Cadre” (2006); Outstanding Scholarship (2006, 2007)

\end{rSection}

%--------------------------------------------------------------------------------
%    Projects And Seminars
%-----------------------------------------------------------------------------------------------
\begin{rSection}{\faBriefcase~PROFESSIONAL EXPERIENCE}
  {\bf Department of Engineering, Indiana University Bloomington} \hfill {Bloomington, IN}
  \\{Incoming Assistant Professor} \hfill {Starting Jan. 2027}

  \vspace{+0.5ex}
  {\bf Department of Computer Science, Saint Louis University (SLU)} \hfill {St. Louis, MO}
  \\{Assistant Professor} \hfill {Aug. 2023--Dec. 2026}

  \vspace{+0.5ex}
  {\bf Innovative Computer Laboratory (ICL), UTK } \hfill {Knoxville, TN}
  \\{Post-Doctoral Research Associate} \hfill {Mar. 2023--Jul. 2023}
  %\\{Advisor: Dr. George Bosilca}
  
  \vspace{+0.5ex}
  {\bf Cerebras Systems, Inc. } \hfill {Sunnyvale, CA}
  \\{Member of Technical Staff} \hfill {Aug. 2022--Jan. 2023}
  
  %\vspace{+0.5ex}
  %{\bf Innovative Computer Laboratory (ICL), UTK } \hfill {Knoxville, TN}
  %\\{Graduate Research Assistant} \hfill {2017 - 2022}

  \vspace{+0.5ex}
  {\bf Cerebras Systems, Inc. } \hfill {Sunnyvale, CA}
  \\{Summer Intern} \hfill {May 2021--Aug. 2021}

  \vspace{+0.5ex}
  {\bf Cadence Design Systems, Inc. } \hfill {Austin, TX}
  \\{Summer Intern} \hfill {May 2020--Jul. 2020}

  %\vspace{+0.5ex}
  %{\bf National University of Defense Technology (NUDT) } \hfill {Changsha, China}
  %\\{HPC Software Developer \& Research Scientist} \hfill {2010 - 2013}
  %\\{HPC Engineer} \hfill {2010 - 2013}

  %\vspace{+0.5ex}
  %{\bf Dongjia Orthopedics} \hfill {Shandong, China}
  %\\{Software Engineer} \hfill {2009 - 2010}
  %\\{Software Engineer} \hfill {Aug. 2009 - Apr. 2010}

\end{rSection}
%----------------------------------------------------------------------------------------
%	TECHNICAL STRENGTHS SECTION
%----------------------------------------------------------------------------------------

%\newpage

\begin{rSection}{\faTrophy ~AWARDS \& GRANTS}

  \begin{itemize}
  \setlength{\itemindent}{-.2in}

  \item {\color{cyan} {\bf AWARDS}}
    \begin{itemize}
      %\setlength\itemsep{-0.5ex}
      \setlength{\itemindent}{-.2in}
      %\item[\ding{43}] 19,440 Node Hours on Frontier Supercomputer, Oak Ridge National Laboratory, US \hfill 2024
      %\item[\ding{43}] 5,000 Node Hours on Polaris Supercomputer, Argonne National Laboratory, US \hfill 2025 
      %\item[\ding{43}] Best Paper Runner-up Award, Workshop on Distributed Runtime Systems for Scientific Big Data \hfill 2025
      \item[\ding{43}] Best Paper Runner-up Award, DRBSD \hfill 2025
      \item[\ding{43}] TCSC SCALE Challenge Award (Finalist, 2 Entries) \hfill 2025
      \item[\ding{43}] {\color{red} ACM Gordon Bell Prize for Climate Modelling} \hfill 2024
      \item[\ding{43}] {\color{red} ACM Gordon Bell Prize (Finalist)} \hfill 2024
      %\item[\ding{43}] {Subaward Funding, Oak Ridge National Laboratory, US \hfill 2024}
      %\item[\ding{43}] 20,000 Node Hours on Frontier Supercomputer, Oak Ridge National Laboratory, US \hfill 2024
      %\item[\ding{43}] 206,780 Node Hours on Summit Supercomputer, Oak Ridge National Laboratory, US \hfill 2024
      %\item[\ding{43}] 20,000 Node Hours on Frontier Supercomputer, Oak Ridge National Laboratory, US \hfill 2024
      %\item[\ding{43}] 72000 Node Hours on Leonardo Supercomputer, EuroHPC/CINECA, Italy \hfill 2024
      %\item[\ding{43}] 3,000,000 Node Hours on Shaheen II Supercomputer, KAUST, Saudi Arabia \hfill 2019 - 2023
      \item[\ding{43}] {\color{red} ACM Gordon Bell Prize (Finalist)} \hfill 2022
      %\item[\ding{43}] 4,000,000 Node Hours on Fugaku Supercomputer, RIKEN, Japan \hfill 2022
      \item[\ding{43}] SIAM Student Travel Award \hfill 2021
      %\item[\ding{43}] 40,000 Node Hours on Summit Supercomputer, Oak Ridge National Laboratory, US \hfill 2021
      \item[\ding{43}] Best Paper Award, CLUSTER, IEEE \hfill 2020
      \item[\ding{43}] Graduate Student Senate (GSS) Travel Awards, UTK \hfill 2020
      %\item[\ding{43}] Honor of Outstanding Graduates, OUC \hfill 2014
      %\item[\ding{43}] Graduate Student Scholarship, OUC \hfill 2014
      %\item[\ding{43}] Honor of Annual Advanced Worker, NUDT \hfill 2010, 2011
      %\item[\ding{43}] Honor of Bronze Medal of TianHe-1A, NUDT \hfill 2010
      %\item[\ding{43}] Outstanding Scholarship, HNU \hfill 2006, 2007
      %\item[\ding{43}] Honor of Excellent Student Cadre, HNU \hfill 2006
    \end{itemize}
    %\hfill
    %\setlength{\itemindent}{-.2in}
    \item {\color{cyan} {\bf GRANTS}}
      \begin{itemize}
        %\setlength\itemsep{-0.5ex}
        \setlength{\itemindent}{-.2in}
       \item[\ding{43}] POSE/NSF, Site-PI \hfill \href{https://www.nsf.gov/awardsearch/show-award?AWD_ID=2449220}{2025-2027}
       \item[\ding{43}] CRII/OAC/NSF, Solo-PI \hfill \href{https://www.nsf.gov/awardsearch/show-award?AWD_ID=2451577}{2025-2027}
       \item[\ding{43}] E3SM/ORNL/DOE, Sub-Award \hfill 2024-2025
       %\item[\ding{43}] PaRSEC/ECP/DOE, Other Personnel \hfill 2017-2022
      \end{itemize}
    %\hfill
    \item {\color{cyan} {\bf COMPUTING RESOURCE ALLOCATIONS}}
    \begin{itemize}
      %\setlength\itemsep{-0.5ex}
      \setlength{\itemindent}{-.2in}
    \item[\ding{43}] Frontier Supercomputer, Oak Ridge National Laboratory, US \hfill 2022-Present
    \item[\ding{43}] Fugaku Supercomputer, RIKEN, Japan  \hfill 2021-Present
    \item[\ding{43}] Polaris Supercomputer, Argonne National Laboratory, US \hfill 2025-Present
    \item[\ding{43}] Leonardo Supercomputer, EuroHPC/CINECA, Italy \hfill 2024
    \item[\ding{43}] Alps Supercomputer, Swiss National Supercomputing Centre, Switzerland \hfill 2024
    \item[\ding{43}] Summit Supercomputer, Oak Ridge National Laboratory, US \hfill 2020-2024
    \end{itemize}
  \end{itemize}
  
\end{rSection}

%\nobibliography{ref}
%\bibliographystyle{unsrt}


\newcommand{\qinglei}{{\bf Qinglei Cao}}

\begin{rSection}{\faBook~Publications}

  \begin{itemize}%[labelindent=0pt, align=left]
    \setlength{\itemindent}{-.2in}
    \item {\color{cyan} \bf Conference Papers}
  
  \begin{enumerate}[leftmargin=8pt, align=left]%[labelindent=-1in, label=\arabic*, align=left]
    %[leftmargin=\parindent,align=left,labelwidth=\parindent,labelsep=2pt]%[leftmargin=-0.5pt]
    %\vspace{-1ex}
    %\setlength\itemsep{-0.2ex}
    %\setlength{\itemindent}{-.2in}
    %
    \item Dawei Gao, Dali Wang, Zhuowei Gu, Qinglei Cao, Peter E. Thornton, Daniel Ricciuto, and Yunhe Feng. HELM: Restart-Compatible Trustworthy AI Emulation for Accelerating Long-Horizon Global Land Simulations. IEEE International Conference on Big Data, 2026
    %
    \item Zhuowei Gu, Aurelien Bouteiller, Joseph Schuchart, Thomas Herault and Qinglei Cao. Merging Parameterized Task Graphs in PaRSEC Through JDF Composition with LLM Assistance. Workshop on Asynchronous Many-Task Systems and Applications(WAMTA), 2026 \hfill[\href{https://qingleicao.com/papers/2026-wamta-Merging Parameterized Task Graphs in PaRSEC Through JDF Composition with LLM Assistance.pdf}{PDF}]
    %
    \item Qiao Zhang, Rabab Alomairy, Dali Wang, Addison Thurston and Qinglei Cao. Leveraging PaRSEC for Task-Based Heterogeneous Computing in Python and Julia. Workshop on Asynchronous Many-Task Systems and Applications(WAMTA), 2026 \hfill[\href{https://qingleicao.com/papers/2026-wamta-Leveraging PaRSEC for Task-Based Heterogeneous Computing in Python and Julia.pdf}{PDF}]
    %
    \item Cory Gardner, Seyun Jeong, Oam Khatavkar, Aiden Moon, Qinglei Cao, and Tae-Hyuk Ahn. Evaluating Accuracy and Performance Tradeoffs in GPU Accelerated Single Cell RNA-seq Analysis. International Workshop on Distributed Runtime Systems for Scientific Big Data (DRBSD, {\bf Best Paper Runner-up Award}), 2025 \hfill[\href{https://qingleicao.com/papers/Evaluating Accuracy and Performance Tradeoffs in GPU Accelerated Single Cell RNA-seq Analysis.pdf}{PDF}]
    %
    \item Rabab Alomairy, Sameh Abdulah, Qinglei Cao, David Keyes and Hatem Ltaief. Sustainably Modeling a Sustainable Future Climate. IEEE High Performance Extreme Computing Conference, 2025 \hfill[\href{https://qingleicao.com/papers/Sustainably Modeling a Sustainable Future Climate.pdf}{PDF}]
    %
    \item Zhuowei Gu, Dali Wang, Dawei Gao, Yunhe Feng and Qinglei Cao. Flexible User-defined Domain Decomposition in Kilometer-Scale E3SM Land Model Simulation. International Conference on Computational Science (ICCS), 2025 \hfill[\href{https://qingleicao.com/papers/Flexible User-defined Domain Decomposition in Kilometer-Scale E3SM Land Model Simulation.pdf}{PDF}]
    %
    \item Rabab Alomairy, Qinglei Cao, Hatem Ltaief and David Keyes. Scalable Low-Rank Solvers for Large-Scale 3D Mesh Deformation Using Global and Compact Support RBF Kernels. IEEE International Symposium on Cluster, Cloud, and Internet Computing (CCGRID, {\bf TCSC SCALE Challenge Award, Finalist}), 2025
    %
    \item Dali Wang, Chen Wang, Qinglei Cao, Jayesh Krishna, Danqing Wu, Weijian Zheng, Peter Schwartz, Fengming Yuan, Kathryn Mohror and Peter Thornton. Scaling Ultrahigh-Resolution E3SM Land Model for Leadership-Class Supercomputers. IEEE International Symposium on Cluster, Cloud, and Internet Computing (CCGRID, {\bf TCSC SCALE Challenge Award, Finalist}), 2025 \hfill[\href{https://qingleicao.com/papers/Scaling Ultrahigh-Resolution E3SM Land Model for Leadership-Class Supercomputers.pdf}{PDF}]
    %
    \item Qiao Zhang, Rabab Alomairy, Dali Wang, Zhuowei Gu and Qinglei Cao. Leveraging Hardware-Aware Computation in Mixed-Precision Matrix Multiply: A Tile-Centric Approach. Workshop on Asynchronous Many-Task Systems and Applications (WAMTA), 2025 \hfill[\href{https://qingleicao.com/papers/2025-wamta-Leveraging Hardware-Aware Computation in Mixed-Precision Matrix Multiply A Tile-Centric Approach.pdf}{PDF}]
    %
    \item Aurelien Bouteiller, Qinglei Cao, Joseph Schuchart and Thomas Herault. Comparing and Contrasting User and Runtime Directed Data Placement Strategies for Owner-Compute, Multi-Accelerator Distributed Task Based Scheduling. Workshop on Asynchronous Many-Task Systems and Applications (WAMTA), 2025 \hfill[\href{https://qingleicao.com/papers/2025-wamta-Comparing and Contrasting User and Runtime Directed Data Placement Strategies for Owner-Compute, Multi-accelerator Distributed Task Based Scheduling.pdf}{PDF}]
    %
    \item Rabab Alomairy, Hatem Ltaief, Qinglei Cao and David Keyes. Analyzing Data Sparsity in Global and Compact Support Radial Basis Functions for 3D Unstructured Mesh Deformation. Workshop on Asynchronous Many-Task Systems and Applications (WAMTA), 2025 \hfill[\href{https://qingleicao.com/papers/2025-wamta-Analyzing Data Sparsity in Global and Compact Support Radial Basis Functions for 3D Unstructured Mesh Deformation.pdf}{PDF}]
    %
    \item Rabab Alomairy, Qinglei Cao, Hatem Ltaief, David E. Keyes, and Alan Edelman. Scalable Hamming Distance Computation Using Accelerated Matrix Transformations. ISC High Performance, 2025 \hfill[\href{https://qingleicao.com/papers/Scalable_Hamming_Distance_Computation_Using_Accelerated_Matrix_Transformations.pdf}{PDF}]
    %
    \item Sameh Abdulah, Allison H. Baker, George Bosilca, Qinglei Cao, Stefano Castruccio, Marc G. Genton, David E. Keyes, Zubair Khalid, Hatem Ltaief, Yan Song, Georgiy L. Stenchikov, and Ying Sun. Boosting Earth System Model Outputs And Saving PetaBytes in their Storage Using Exascale Climate Emulators. International Conference for High Performance Computing, Networking, Storage and Analysis (SC, {\bf ACM Gordon Bell Prize for Climate Modelling}), 2024  \hfill[\href{https://qingleicao.com/papers/Boosting Earth System Model Outputs And Saving PetaBytes in Their Storage Using Exascale Climate Emulators.pdf}{PDF}]
    %
    \item Hatem Ltaief, Rabab Alomairy, Qinglei Cao, Jie Ren, Lotfi Slim, Thorsten Kurth, Benedikt Dorschner, Salim Bougouffa, Rached Abdelkhalek, and David E. Keyes. Toward Capturing Genetic Epistasis From Multivariate Genome-Wide Association Studies Using Mixed-Precision Kernel Ridge Regression. International Conference for High Performance Computing, Networking, Storage and Analysis (SC, {\bf ACM Gordon Bell Prize, Finalist}), 2024 \hfill[\href{https://qingleicao.com/papers/Toward Capturing Genetic Epistasis From Multivariate Genome-Wide Association Studies Using Mixed-Precision Kernel Ridge Regression.pdf}{PDF}]
    %
    \item Shihui Song, Yafan Huang, Peng Jiang, Xiaodong Yu, Weijian Zheng, Sheng Di, Qinglei Cao, Yunhe Feng, Zhen Xie, and Franck Cappello. CereSZ: Enabling and Scaling Error-bounded Lossy Compression on Cerebras CS-2. International Symposium on High-Performance Parallel and Distributed Computing (HPDC), 2024 \hfill[\href{https://qingleicao.com/papers/CereSZ- Enabling and Scaling Error-bounded Lossy Compression on Cerebras CS-2.pdf}{PDF}]
    %
    \item Kareem Shaik, Dali Wang, Weijian Zheng, Qinglei Cao, Heng Fan, Peter Schwartz, and Yunhe Feng. S3LLM : Large-Scale Scientific Software Understanding with LLMs using Source, Metadata, and Document. International Conference on Computational Science (ICCS), 2024 \hfill[\href{https://qingleicao.com/papers/S3LLM- Large-Scale Scientific Software Understanding with LLMs using Source, Metadata, and Document.pdf}{PDF}]
    %
    \item Qinglei Cao, Thomas Herault, Aurelien Bouteiller, Joseph Schuchart, and George Bosilca. Evaluating PaRSEC through Matrix Computations in Scientific Applications. Workshop on Asynchronous Many-Task Systems and Applications (WAMTA), 2024 \hfill[\href{https://qingleicao.com/papers/Evaluating PaRSEC Through Matrix Computations in Scientific Applications.pdf}{PDF}]
    % 
    \item {Qinglei Cao}, Sameh Abdulah, Hatem Ltaief, Marc G Genton, David E Keyes, and George Bosilca. Reducing Data Motion and Energy Consumption of Geospatial Modeling Applications Using Automated Precision Conversion. IEEE International Conference on Cluster Computing ({CLUSTER}), 2023 \hfill[\href{https://qingleicao.com/papers/Reducing Data Motion and Energy Consumption of Geospatial Modeling Applications Using Automated Precision Conversion.pdf}{PDF}]
    %
    \item {Qinglei Cao}, Sameh Abdulah, Rabab Alomairy, Yu Pei, Pratik Nag, George Bosilca, Jack Dongarra, Marc G Genton, David E Keyes, Hatem Ltaief, and Ying Sun. Reshaping geostatistical modeling and prediction for extreme-scale environmental applications. International Conference for High Performance Computing, Networking, Storage and Analysis (SC, {\bf ACM Gordon Bell Prize, Finalist}), 2022 \hfill[\href{https://qingleicao.com/papers/Reshaping geostatistical modeling and prediction for extreme-scale environmental applications.pdf}{PDF}]
    %
    \item {Qinglei Cao}, Rabab Alomairy, Yu Pei, George Bosilca, Hatem Ltaief, David Keyes, and Jack Dongarra. A framework to exploit data sparsity in tile low-rank Cholesky factorization. IEEE International Parallel \& Distributed Processing Symposium ({IPDPS}), 2022 \hfill[\href{https://qingleicao.com/papers/A Framework to Exploit Data Sparsity in Tile Low-Rank Cholesky Factorization.pdf}{PDF}]
    %
    \item {Qinglei Cao}, Yu Pei, Kadir Akbudak, George Bosilca, Hatem Ltaief, David Keyes, and Jack Dongarra. Leveraging parsec runtime support to tackle challenging 3d data-sparse matrix problems. IEEE International Parallel and Distributed Processing Symposium ({IPDPS}), 2021 \hfill[\href{https://qingleicao.com/papers/Leveraging PaRSEC Runtime Support to Tackle Challenging 3D Data-Sparse Matrix Problems.pdf}{PDF}]
    % 
    \item Yunhe Feng, Dong Zhong, Peng Sun, Weijian Zheng, {Qinglei Cao}, Xi Luo, and Zheng Lu. Micromobility in smart cities: A closer look at shared dockless e-scooters via big social data. IEEE International Conference on Communications (ICC), 2021 \hfill[\href{https://qingleicao.com/papers/Micromobility in Smart Cities A Closer Look at Shared Dockless E-Scooters via Big Social Data.pdf}{PDF}]
    %
    \item {Qinglei Cao}, George Bosilca, Wei Wu, Dong Zhong, Aurelien Bouteiller, and Jack Dongarra. Flexible data redistribution in a task-based runtime system. IEEE International Conference on Cluster Computing ({CLUSTER}), 2020 \hfill[\href{https://qingleicao.com/papers/Flexible Data Redistribution in a Task-Based Runtime System.pdf}{PDF}]
    %
    \item Qinglei Cao, Yu Pei, Kadir Akbudak, Aleksandr Mikhalev, George Bosilca, Hatem Ltaief, David Keyes, and Jack Dongarra. Extreme-scale task-based Cholesky factorization toward climate and weather prediction applications. ACM Platform for Advanced Scientific Computing Conference (PASC), 2020 \hfill[\href{https://qingleicao.com/papers/Extreme-Scale Task-Based Cholesky Factorization Toward Climate and Weather Prediction Applications.pdf}{PDF}]
    %
    \item Elliott Slaughter, Wei Wu, Yuankun Fu, Legend Brandenburg, Nicolai Garcia, Wilhem Kautz, Emily Marx, Kaleb S. Morris, {Qinglei Cao}, George Bosilca, Seema Mirchandaney, Wonchan Lee, Sean Treichler, Patrick McCormick, and Alex Aiken. Task bench: a parameterized benchmark for evaluating parallel runtime performance. IEEE/ACM International Conference for High Performance Computing, Networking, Storage and Analysis ({SC}), 2020 \hfill[\href{https://qingleicao.com/papers/Task_Bench_A_Parameterized_Benchmark_for_Evaluating_Parallel_Runtime_Performance.pdf}{PDF}]
    %
    \item Xi Luo, Wei Wu, George Bosilca, Yu Pei, {Qinglei Cao}, Thananon Patinyasakdikul, Dong Zhong, and Jack Dongarra. Han: a hierarchical autotuned collective communication framework. IEEE International Conference on Cluster Computing (CLUSTER, {\bf Best Paper}), 2020 \hfill[\href{https://qingleicao.com/papers/HAN_a_Hierarchical_AutotuNed_Collective_Communication_Framework.pdf}{PDF}]
    %
    \item Dong Zhong, {Qinglei Cao}, George Bosilca, and Jack Dongarra. Using advanced vector extensions AVX-512 for MPI reductions. ACM European MPI Users’ Group Meeting (EuroMPI), 2020 \hfill[\href{https://qingleicao.com/papers/Using advanced vector extensions AVX-512 for MPI reductions.pdf}{PDF}]
    %
    \item Dong Zhong, Pavel Shamis, {Qinglei Cao}, George Bosilca, Shinji Sumimoto, Kenichi Miura, and Jack Dongarra. Using ARM scalable vector extension to optimize OpenMPI. IEEE/ACM International Symposium on Cluster, Cloud and Internet Computing ({CCGRID}), 2020 \hfill[\href{https://qingleicao.com/papers/Using_Arm_Scalable_Vector_Extension_to_Optimize_OPEN_MPI.pdf}{PDF}]
    %
    \item Yu Pei, {Qinglei Cao}, George Bosilca, Piotr Luszczek, Victor Eijkhout, and Jack Dongarra. Communication avoiding 2d stencil implementations over PaRSEC task-based runtime. IEEE International Parallel and Distributed Processing Symposium Workshops (IPDPSW), 2020 \hfill[\href{https://qingleicao.com/papers/Communication_Avoiding_2D_Stencil_Implementations_over_PaRSEC_Task-Based_Runtime.pdf}{PDF}]
    %
    \item {Qinglei Cao}, Yu Pei, Thomas Herault, Kadir Akbudak, Aleksandr Mikhalev, George Bosilca, Hatem Ltaief, David Keyes, and Jack Dongarra. Performance analysis of tile low-rank Cholesky factorization using parsec instrumentation tools. IEEE/ACM International Workshop on Programming and Performance Visualization Tools (ProTools at SC), 2019 \hfill[\href{https://qingleicao.com/papers/Performance Analysis of Tile Low-Rank Cholesky Factorization Using PaRSEC Instrumentation Tools.pdf}{PDF}]
    %
    %\item Yan Yan, Jie Nie, Lei Huang, Zhen Li, {Qinglei Cao}, and Zhiqiang Wei. Exploring relationship between face and trustworthy impression using mid-level facial features. International Conference on Multimedia Modeling (MMM), 2016 \hfill[\href{https://qingleicao.com/papers/Exploring relationship between face and trustworthy impression using mid-level facial features.pdf}{PDF}]
    %
    %\item Qinglei Cao, Tiantian Li, and Bo Yin. "Parallel optimization research based on numerical simulation." In 2016 Third International Conference on Electrical, Electronics, Computer Engineering and their Applications (EECEA), 2016
    %
    %\item Yan Yan, Jie Nie, Lei Huang, Zhen Li, {Qinglei Cao}, and Zhiqiang Wei. Is your first impression reliable? trustworthy analysis using facial traits in portraits. International Conference on Multimedia Modeling (MMM), 2015 \hfill[\href{https://qingleicao.com/papers/Is your first impression reliable.pdf}{PDF}]
    %
    %\item Zhenyu Liu, {Qinglei Cao}, Yang Pan, and Gui Chen. "Security and Privacy Protection of Smart Home Based on IPv6." In 2014 International Conference on Mechatronics, Electronic, Industrial and Control Engineering (MEIC), 2014
  \end{enumerate}

  %\item {\color{cyan} \bf Journal Papers}

  %\begin{enumerate}[leftmargin=8pt, align=left]%[labelindent=10pt, label=\arabic*, align=left]
    %\item Aurelien Bouteiller, Thomas Herault, Qinglei Cao, Joseph Schuchart, and George Bosilca. PaRSEC: Scalability, Flexibility, and Hybrid Architecture Support for Task-based Applications in ECP. International Journal of High Performance Computing Applications (IJHPCA), 2024
    %
    %\item {Qinglei Cao}, George Bosilca, Nuria Losada, Wei Wu, Dong Zhong, and Jack Dongarra. Evaluating data redistribution in parsec. IEEE Transactions on Parallel and Distributed Systems ({TPDS}), 2022\hfill[\href{https://qingleicao.com/papers/Evaluating Data Redistribution in PaRSEC.pdf}{PDF}]
    %
    %\item Sameh Abdulah, {Qinglei Cao}, Yu Pei, George Bosilca, Jack Dongarra, Marc G. Genton, David E. Keyes, Hatem Ltaief, and Ying Sun. Accelerating geostatistical modeling and prediction with mixed-precision computations: A high-productivity approach with parsec. IEEE Transactions on Parallel and Distributed Systems ({TPDS}), 2022 \hfill[\href{https://qingleicao.com/papers/Accelerating Geostatistical Modeling and Prediction With Mixed-Precision Computations- A High-Productivity Approach with PaRSEC.pdf}{PDF}]
    %
    %\item Dong Zhong, {Qinglei Cao}, George Bosilca, and Jack Dongarra. Using long vector extensions for MPI reductions. Parallel Computing ({PARCO}), 2021 \hfill[\href{https://qingleicao.com/papers/Using long vector extensions for MPI reductions.pdf}{PDF}]
  %R\end{enumerate}

  \item {\color{cyan} \bf Journal Papers}
  
  \begin{enumerate}[leftmargin=8pt, align=left]%[labelindent=-1in, label=\arabic*, align=left]
    %[leftmargin=\parindent,align=left,labelwidth=\parindent,labelsep=2pt]%[leftmargin=-0.5pt]
    %\vspace{-1ex}
    %\setlength\itemsep{-0.2ex}
    %\setlength{\itemindent}{-.2in}
    %
    \item Dali Wang, Chen Wang, Qinglei Cao, Peter Schwartz, Fengming Yuan, Daniel Ricciuto, Peter Thornton, Shih-Chieh Kao, Jayesh Krishna, and Weijian Zheng. Enabling Kilometer-Scale E3SM Land Model Simulation over North America: A New Integrated Framework Solution. International Journal of High Performance Computing Applications (IJHPCA), 2026 \hfill[\href{https://qingleicao.com/papers/Enabling Kilometer-Scale E3SM Land Model Simulation over North America A New Integrated Framework Solution.pdf}{PDF}]
    %
    \item Dali Wang, Qian Gong, Zirui Liu, Xiao Wang, Qinglei Cao and Scott Klasky. Scalable foundation models for numerical simulations on HPC platforms. Frontiers, 2026 \hfill[\href{https://qingleicao.com/papers/Scalable foundation models for numerical simulations on HPC platforms.pdf}{PDF}]
    %
    \item Qiao Zhang, Rabab Alomairy, Dali Wang, Zhuowei Gu, and Qinglei Cao. High-performance mixed-precision matrix multiplication via tile-centric design on modern architectures. SN Computer Science. Springer, 2026 \hfill[\href{https://qingleicao.com/papers/High-Performance Mixed-Precision Matrix Multiplication via TileCentric Design on Modern Architectures.pdf}{PDF}]
    %
    \item Aurelien Bouteiller, Thomas Herault, Qinglei Cao, Joseph Schuchart, and George Bosilca. PaRSEC: Scalability, Flexibility, and Hybrid Architecture Support for Task-based Applications in ECP. International Journal of High Performance Computing Applications (IJHPCA), 2025 \hfill[\href{https://dl.acm.org/doi/10.1177/10943420241290520}{PDF}]
    %
    \item Wang, Dali, Peter Schwartz, Fengming Yuan, Franklin Eaglebarge, Danial Riccuito, Peter Thornton, Chris Layton, and Qinglei Cao. Portable Software Environment for Ultrahigh-Resolution ELM Development on GPUs. Journal of Computer and Communications, 2025 \hfill[\href{https://www.scirp.org/journal/paperinformation?paperid=140718}{PDF}]
    %
    \item Sameh Abdulah, Qinglei Cao, Yu Pei, George Bosilca, Jack Dongarra, Marc G. Genton, David E. Keyes, Hatem Ltaief, and Ying Sun. Accelerating geostatistical modeling and prediction with mixed-precision computations: A high-productivity approach with parsec. IEEE Transactions on Parallel and Distributed Systems ({TPDS}), 2022 \hfill[\href{https://qingleicao.com/papers/Accelerating Geostatistical Modeling and Prediction With Mixed-Precision Computations- A High-Productivity Approach with PaRSEC.pdf}{PDF}]
    %
    \item Dong Zhong, {Qinglei Cao}, George Bosilca, and Jack Dongarra. Using long vector extensions for MPI reductions. Parallel Computing ({PARCO}), 2021 \hfill[\href{https://qingleicao.com/papers/Using long vector extensions for MPI reductions.pdf}{PDF}]
    %
   \item {Qinglei Cao}, George Bosilca, Nuria Losada, Wei Wu, Dong Zhong, and Jack Dongarra. Evaluating data redistribution in parsec. IEEE Transactions on Parallel and Distributed Systems ({TPDS}), 2021\hfill[\href{https://qingleicao.com/papers/Evaluating Data Redistribution in PaRSEC.pdf}{PDF}]
   %
  \end{enumerate}


  \item {\color{cyan} \bf Abstracts \& Posters} 
  
  \begin{enumerate}[leftmargin=8pt, align=left]%[labelindent=10pt, label=\arabic*, align=left]
    %
    \item Dali Wang, Peter E. Thornton, Daniel M. Ricciuto, Xiaojuan Yang, Xiao Wang, Dan Lu, Dawei Gao, Yunhe Feng, Qinglei Cao, Zhuowei Gu, Zirui Liu. Accelerating Earth System Simulations: Creating an AI-Powered Emulator for the E3SM Land Model. American Geophysical Union (AGU), 2025
    %
    \item Dali Wang, Peter Schwartz, Fengming Yuan, Daniel M. Ricciuto, Shih-Chieh Kao, Michele Thornton, Anthony Walker, Peter E. Thornton, Qinglei Cao, and Chen Wang. Kilometer-scale E3SM Land Model: Code Development, Deployment, Evaluation, and Applications. American Geophysical Union (AGU), 2024
    %
    %\item Qinglei Cao, Dali Wang, Zhuowei Gu, Yunhe Feng, Weijian Zheng, and Peter Schwart. Mixed-Precision for Extreme-Scale Geostatistical Modeling on Heterogeneous Platforms. International Congress on Environmental Modelling and Software (iEMSs), 2024
    %
    %\item Dali Wang, Fengming Yuan, Peter Schwartz, Daniel Ricciuto, Peter Thornton, Shih-Chieh Kao, Michele Thornton, Anthony Walker, Qinglei Cao, and Zhuowei Gu. Enable Kilometre-scale E3SM Land Simulation with User-defined Datasets and KiloCraft. International Congress on Environmental Modelling and Software (iEMSs), 2024
    %
    \item {Qinglei Cao}, Yu Pei, Kadir Akbudak, Aleksandr Mikhalev, George Bosilca, Hatem Ltaief, David Keyes, and Jack Dongarra. Extreme-Scale Tile Low-Rank Cholesky Factorization Using the PaRSEC Task-Based Runtime. ACM Platform for Advanced Scientific Computing Conference ({PASC}), 2021
  \end{enumerate}

\end{itemize}

\end{rSection}


\begin{rSection}{\faHourglass[1] ~PROFESSIONAL ACTIVITIES}
  
  \begin{itemize}
    %\setlength\itemsep{-0.5ex}
    \setlength{\itemindent}{-.2in}
    \item {\color{cyan} {\bf Editorial Board/Guest Editor}}
      \begin{itemize}
        %\setlength\itemsep{-0.5ex}
        \setlength{\itemindent}{-.2in}
       \item[\ding{43}] SN Computer Science \hfill 2025 - 2026
       \item[\ding{43}] Electronics, MDPI \hfill 2026
       \item[\ding{43}] American Journal of Computer Science and Technology \hfill 2023 - 2024
       \end{itemize}
    %
    %\hfill
    \item {\color{cyan} {\bf Organizer}}
      \begin{itemize}
        %\setlength\itemsep{-0.5ex}
        \setlength{\itemindent}{-.2in}
       \item[\ding{43}] Workshop on Fostering an Open-Source Runtime Eco-System for PaRSEC \hfill 2026
       \item[\ding{43}] Workshop on Asynchronous Many-Task Systems and Applications \hfill 2025
       \end{itemize}
    %
    %\hfill
    \item {\color{cyan} {\bf Technical Program Chair}}
      \begin{itemize}
        %\setlength\itemsep{-0.5ex}
        \setlength{\itemindent}{-.2in}
       \item[\ding{43}] Workshop on Accelerator Programming and Directives (WACCPD at SC) \hfill 2026
       \item[\ding{43}] Workshop on Asynchronous Many-Task Systems and Applications \hfill 2025
       \end{itemize}
    %
    %\hfill
    \item {\color{cyan} {\bf Technical Program Committee}}
      \begin{itemize}
        %\setlength\itemsep{-0.5ex}
        \setlength{\itemindent}{-.2in}
        \item[\ding{43}] International Conference on Supercomputing (ICS) \hfill 2027
        \item[\ding{43}] International Parallel \& Distributed Processing Symposium (IPDPS) \hfill 2024, 2027
        \item[\ding{43}] SuperComputing Asia (SCA) \hfill 2027
        \item[\ding{43}] SuperComputing India (SCI) \hfill 2027
        \item[\ding{43}] Workshop on Asynchronous Many-Task Systems and Applications \hfill 2024, 2026, 2027
        \item[\ding{43}] Intl Conference for High Performance Computing, Networking, Storage and Analysis (SC) \hfill 2025, 2026
        \item[\ding{43}] International Supercomputing Conference (ISC) High Performance \hfill 2024, 2026
        \item[\ding{43}] Workshop on HPC on Heterogeneous Hardware \hfill 2022 - 2026
        \item[\ding{43}] Twelfth Workshop on Accelerator Programming and Directives (WACCPD)\hfill 2025
        \item[\ding{43}] International Conference on Parallel Processing (ICPP) \hfill 2024, 2025
        \item[\ding{43}] International Conference on Cluster Computing (CLUSTER, Poster) \hfill 2025
        \item[\ding{43}] Asynchronous Many-Task systems for Exascale (AMTE) \hfill 2025
        \item[\ding{43}] Association for the Advancement of Artificial Intelligence (AAAI) Undergraduate Consortium \hfill 2024
        \item[\ding{43}] Intl Conference for High Performance Computing, Networking, Storage and Analysis AD/AE ({SC}) \hfill 2021
        \item[\ding{43}] Intl Conference on Advances and Trends in Software Engineering (SOFTENG) \hfill 2021 - 2024
      \end{itemize}
    %
    %\hfill
    \item {\color{cyan} {\bf Conference \& Journal Reviewer}}
      \begin{itemize}
        %\setlength\itemsep{-0.5ex}
        \setlength{\itemindent}{-.2in}
        %\item[\ding{43}] Computational and Structural Biotechnology Journal \hfill 2024
        \item[\ding{43}] ACM Web Conference \hfill 2027
        \item[\ding{43}] ACM Transactions on Architecture and Code Optimization \hfill 2026
        \item[\ding{43}] IEEE Transactions on Parallel and Distributed Systems \hfill 2026
        \item[\ding{43}] ACM Transactions on Mathematical Software (TOMS) \hfill 2023, 2026
        \item[\ding{43}] Parallel Computing \hfill 2023
        \item[\ding{43}] Journal of Supercomputing \hfill 2023
        \item[\ding{43}] IEEE Transactions on Multimedia \hfill 2023
        \item[\ding{43}] International Conference on Emerging Information Security and Applications \hfill 2022
        \item[\ding{43}] Intl Conference for High Performance Computing, Networking, Storage and Analysis ({SC}) \hfill 2020, 2021
        \item[\ding{43}] PeerJ Computer Science \hfill 2021
        \item[\ding{43}] International Conference on Cluster Computing ({CLUSTER}) \hfill 2020
        \item[\ding{43}] International Conferences on High Performance Computing and Communications (HPCC) \hfill 2020, 2021
        \end{itemize}      
    %
    %\hfill
    \item {\color{cyan} {\bf Miscellaneous}}
      \begin{itemize}
        %\setlength\itemsep{-0.5ex}
        \setlength{\itemindent}{-.2in}
       %\item[\ding{43}] {\color{red} Application Support for Student Cluster Competition (SCC) and IndySCC} at SC\hfill 2025
       \item[\ding{43}] {\bf Application Support for Student Cluster Competition (SCC) and IndySCC} at SC\hfill 2025
       \item[\ding{43}] Section Chair at SC\hfill 2025
       \end{itemize}
    
  \end{itemize}
  
  
\end{rSection}



\begin{rSection}{\faDiamond ~TEACHING EXPERIENCE}
  
  \begin{itemize}%[leftmargin=-0.5pt]
    %\vspace{-1ex}
    %\setlength\itemsep{-1.1ex}
    \setlength{\itemindent}{-.2in}
    \item {\color{cyan} {\bf Lecturer}}
    \begin{itemize}
      \setlength{\itemindent}{-.2in}
      \item[\ding{43}] CSCI 2510 {Principles of Computing Systems} \hfill Spring (2025, 2026), Fall (2025, 2026), SLU
      \item[\ding{43}] CSCI 4620/5620 {Distributed Computing} \hfill Spring (2024), Fall (2023 - 2026), SLU
      %\item[\ding{43}] CSCI 2510 {Principles of Computing Systems} \hfill Spring 2025, SLU
      %\item[\ding{43}] CSCI 4620/5620 {Distributed Computing} \hfill Fall 2024, SLU
      %\item[\ding{43}] CSCI 4620/5620 {Distributed Computing} \hfill Spring 2024, SLU
      %\item[\ding{43}] CSCI 4620/5620 {Distributed Computing} \hfill Fall 2023, SLU
    \end{itemize}    
    %
    \item {\color{cyan} {\bf Guest Lecturer}}
    \begin{itemize}
      %\setlength\itemsep{-1.1ex}
      \setlength{\itemindent}{-.2in}
      \item[\ding{43}] CSCE 5300 {Introduction to Big Data and Data Science} \hfill Spring (2023, 2025), Fall (2023), UNT
      \item[\ding{43}] CSCI 5090 {Computer Science Colloquium} \hfill Fall (2023, 2024), SLU 
      %\item[\ding{43}] CSCI 5090 {Computer Science Colloquium} \hfill Fall 2023, SLU 
      %\item[\ding{43}] CSCE 5300 {Introduction to Big Data and Data Science} \hfill Fall 2023, UNT
      %\item[\ding{43}] CSCE 5300 {Introduction to Big Data and Data Science} \hfill Spring 2023, UNT
    \end{itemize}
    %
    \item {\color{cyan} {\bf Teaching Assistant}}
    \begin{itemize}
      %\setlength\itemsep{-1.1ex}
      \setlength{\itemindent}{-.2in}
      \item[\ding{43}] COSC 594 {Scientific Computing for Engineers} \hfill Spring 2018, UTK
      \item[\ding{43}] COSC 361 {Operating Systems} \hfill Spring (2017), Fall (2016), UTK
      %\item[\ding{43}] COSC 361 {Operating Systems} \hfill Fall 2016, UTK
    \end{itemize}
    %
  \end{itemize}   
  
\end{rSection}

%\newpage

\begin{rSection}{\faConnectdevelop~PRESENTATIONS \& TALKS}
  
  \begin{itemize}%align=left]%[leftmargin=-0.5pt]
    %\vspace{-1ex}
    %\setlength\itemsep{-1.1ex}
    \setlength{\itemindent}{-.2in}
    %
    \item {\color{cyan} {\bf Talk}}
    \begin{itemize}
      %\setlength\itemsep{-1.1ex}
      \setlength{\itemindent}{-.2in}
      %\item[\ding{43}] {IndySCC and SCC Webinar at SC'25\hfill 2025}
      \item[\ding{43}] {Stevens Institute of Technology\hfill 2025}
      \item[\ding{43}] {OLCF User Meeting\hfill 2025}
      \item[\ding{43}] {ICL 35th Anniversary Workshop  \hfill 2025}
      \item[\ding{43}] {NVIDIA GTC (GPU Technology Conference)  \hfill 2025}
      \item[\ding{43}] {SIAM Conference on Computational Science and Engineering \hfill 2025}
      \item[\ding{43}] {University of North Texas \hfill 2025}
      \item[\ding{43}] {Annual Workshop on Charm++ and Its Applications \hfill 2024}
      \item[\ding{43}] {Innovative Computer Laboratory (ICL) Lunch Talk \hfill 2019, 2020, 2021, 2022}
      \item[\ding{43}] Joint Laboratory on Extreme Scale Computing Workshop ({JLESC}) \hfill 2021
      \item[\ding{43}] SIAM Conference on Computational Science and Engineering ({CSE}) \hfill 2021
      \item[\ding{43}] SIAM Conference on Parallel Processing for Scientific Computing (PP) \hfill 2020
    \end{itemize} 
    %
    %\hfill
    \item {\color{cyan} {\bf Paper Presentation}}
    \begin{itemize}
      %\setlength\itemsep{-1.1ex}
      \setlength{\itemindent}{-.2in}
      \item[\ding{43}] {Workshop on Asynchronous Many-Task Systems and Applications ({WAMTA}) \hfill 2024}
      \item[\ding{43}] {International Conference on Cluster Computing ({CLUSTER}) \hfill 2023}
      \item[\ding{43}] {International Parallel and Distributed Processing Symposium ({IPDPS}) \hfill 2021, 2022}
      \item[\ding{43}] {International Conference on Cluster Computing ({CLUSTER}) \hfill 2020}
      \item[\ding{43}] {Platform for Advanced Scientific Computing Conference ({PASC}) \hfill 2020}
      \item[\ding{43}] {International Workshop on Programming and Performance Visualization Tools (ProTools at SC) \hfill 2019}
    \end{itemize}
    %
    %\hfill
    \item {\color{cyan} {\bf Poster}}
    \begin{itemize}
      %\setlength\itemsep{-1.1ex}
      \setlength{\itemindent}{-.2in}
      \item[\ding{43}] Joint Laboratory on Extreme Scale Computing Workshop ({JLESC}) \hfill 2020
      \item[\ding{43}] Platform for Advanced Scientific Computing Conference ({PASC}) \hfill 2020
    \end{itemize}       
  \end{itemize}  

\end{rSection}


\begin{rSection}{\faHeart~Open Source Contributions}
  
  \begin{itemize}%[align=left]
    %\vspace{-1ex}
    %\setlength\itemsep{-1.1ex}
    \setlength{\itemindent}{-.2in}
    \item {\bf [\href{https://github.com/ICLDisco/parsec}{PaRSEC}]}: Task-based runtime system, funded by Exascale Computing Project ({ECP}) 
    \item {\bf [\href{https://github.com/ICLDisco/dplasma}{DPLASMA}]}: Leading implementation of a dense linear algebra package for distributed system
    \item {\bf [\href{https://github.com/ecrc/hicma-x}{HiCMA}]}: Hierarchical math library of exploiting the data sparsity of the matrix operator
    \item {\bf [\href{https://github.com/ecrc/exageostat}{ExaGeostat}]}: Parallel high performance unified framework for computational geostatistics% on many-core systems
  \end{itemize}  
  
\end{rSection}

\begin{comment}
\begin{rSection}{\faKey  ~SKILLS \& PROFICIENCY}

  \begin{itemize}%[align=left]
    %\setlength\itemsep{-1.1ex}
    \setlength{\itemindent}{-.2in}
    \item {\bf Languages:} C (major language), Python, MPI, CUDA, Pthread, OpenMP, C++, Shell
    \item {\bf Software \& Tools:} Linux, Git, Cmake, Makefile, Matlab, GDB, VTune
  \end{itemize}

\end{rSection}


\begin{rSection}{\faKey  HPC RESOURCES}

  \begin{itemize}%[align=left]
    %\setlength\itemsep{-1.1ex}
    \setlength{\itemindent}{-.2in}
    \item {\bf Frontier:} 
    \item {\bf Fugaku:}
    \item {\bf Summit:}
    \item {\bf Shaheen II:}
  \end{itemize}

\end{rSection}

\end{comment}


\begin{rSection}{\faNewspaperO ~MEDIA COVERAGE}
  
  \begin{itemize}%[align=left]
    %\vspace{-1ex}
    %\setlength\itemsep{-1.1ex}
    \setlength{\itemindent}{-.2in}
    \item SC24 Gordon Bell Prizes Awarded\hfill[\href{https://www.hpcwire.com/2024/12/03/sc24-gordon-bell-prizes-awarded/}{HPCwire}]
    %\item NSF NCAR researcher part of team that wins the 2024 ACM Gordon Bell Prize for Climate Modeling\hfill[\href{https://www.cisl.ucar.edu/news/nsf-ncar-researcher-part-team-wins-2024-acm-gordon-bell-prize-climate-modeling}{NCAR}]
    %\item KAUST awarded Gordon Bell Prize for Climate Modelling\hfill[\href{https://www.middleeastainews.com/p/kaust-awarded-gordon-bell-prize}{Middle East AI News}]
    %\item Professional Notes: November 2024\hfill[\href{https://www.slu.edu/news/announcements/2024/november/professional-notes-november-2024.php}{SLU News}]
    \item 2024 Gordon Bell Prizes Awarded at SC24\hfill[\href{https://www.ifipnews.org/2024-gordon-bell-prizes-awarded-at-sc24-2/}{ACM}]
    \item Gordon Bell Prize for Climate Modelling Goes to Team for Exascale Emulator Breakthrough\hfill[\href{https://www.hpcwire.com/off-the-wire/gordon-bell-prize-for-climate-modelling-goes-to-team-for-exascale-emulator-breakthrough/}{HPCwire}]
    %\item Gordon Bell Climate Prize Goes to KAUST Frontier Users’ Exascale Climate Emulator\hfill[\href{https://www.ornl.gov/news/gordon-bell-climate-prize-goes-kaust-frontier-users-exascale-climate-emulator#:~:text=The%202024%20Gordon%20Bell%20Prize,resolution%20but%20without%20the%20computational}{ORNL}]
    \item Recipients of Prestigious Climate Modelling Prize Developed a Technique to Provide More Accurate and Detailed Climate Change Predictions\hfill[\href{https://www.acm.org/media-center/2024/november/gordon-bell-climate-2024}{ACM}]
    \item ACM Gordon Bell Prize for Climate Modelling\hfill[\href{https://awards.acm.org/bell-climate}{ACM}]
    \item ACM Presents Winners of Gordon Bell Climate Modelling Prize\hfill[\href{https://insidehpc.com/2024/11/acm-presents-winners-of-gordon-bell-climate-modelling-prize/}{insideHPC}]
    \item ORNL’s Frontier Powers KAUST-Led Genome Study for Gordon Bell Prize Nomination\hfill[\href{https://www.hpcwire.com/off-the-wire/ornls-frontier-powers-kaust-led-genome-study-for-gordon-bell-prize-nomination/}{HPCwire}]
    \item Bigger, Faster, Smarter Genetics Research\hfill[\href{https://www.olcf.ornl.gov/2024/11/01/bigger-faster-smarter-genetics-research/}{ORNL}]
    \item Frontier Users’ Exascale Climate Emulator Nominated for Gordon Bell Climate Prize\hfill[\href{https://www.olcf.ornl.gov/2024/10/30/frontier-users-exascale-climate-emulator-nominated-for-gordon-bell-climate-prize/}{ORNL}]
    %\item Presenting the Finalists for the 2024 Gordon Bell Prize\hfill[\href{https://sc24.supercomputing.org/2024/10/presenting-the-finalists-for-the-2024-gordon-bell-prize/}{SC24}]
    %\item Three Finalists for the 2024 Gordon Bell Prize for Climate Modelling\hfill[\href{https://sc24.supercomputing.org/2024/10/three-finalists-for-the-2024-gordon-bell-prize-for-climate-modelling/}{SC24}]
    \item Researchers Benchmark Nvidia’s GH200 Supercomputing Chips\hfill[\href{https://www.hpcwire.com/2024/09/04/researchers-benchmark-nvidias-gh200-supercomputing-chips/}{HPCwire}]
    %\item Gordon Bell Award finalists develop method for more efficient computing\hfill[\href{https://www.eurekalert.org/news-releases/970914}{AAAS Eurekalert}]
    \item Gordon Bell Prize Finalists Develop Method for More Efficient Computing\hfill[\href{https://www.eurekalert.org/news-releases/970914}{AAAS Eurekalert}][\href{https://www.hlrs.de/news/detail/gordon-bell-award-finalists-develop-method-for-more-efficient-computing}{HLRS News}]
    %\item Die Finalisten des Gordon Bell Award entwickeln eine Methode für\hfill[\href{https://www.gamingdeputy.com/gr/die-finalisten-des-gordon-bell-award-entwickeln-eine-methode-fuer/}{Gamingdeputy Germany}]
    %\item KAUST Supercomputing Expertise Shines at SC22\hfill[\href{https://www.kaust.edu.sa/en/news/kaust-supercomputing-expertise-shines-at-sc22}{KAUST News}]
    %\item HPE to Build 100+ Petaflops Shaheen III Supercomputer\hfill[\href{https://www.hpcwire.com/2022/09/27/hpe-to-build-100-petaflops-shaheen-iii-supercomputer/}{HPCwire}]
    \item Inside the Gordon Bell Prize Finalist Projects\hfill[\href{https://www.hpcwire.com/2022/09/07/inside-the-gordon-bell-prize-finalist-projects/}{HPCwire}]
    %\item {\begin{figure}[H] \includegraphics[scale=0.25]{figures/riken_news.jpg}\end{figure}}\hfill[\href{https://www.r-ccs.riken.jp/outreach/topics/20220826-1/}{RIKEN}]
    \item SC22 Unveils ACM Gordon Bell Prize Finalists\hfill[\href{https://www.hpcwire.com/2022/08/12/sc22-unveils-acm-gordon-bell-prize-finalists/}{HPCwire}]
    \item 2022 ACM Gordon Bell Prize Finalists Announced\hfill[\href{https://cacm.acm.org/careers/263693-2022-acm-gordon-bell-prize-finalists-announced/fulltext}{Communications of the ACM}]
    \item What’s New in HPC Research: EXA2PRO, DQRA, and HiCMA-PaRSE Frameworks \& More\hfill[\href{https://www.hpcwire.com/2022/06/28/whats-new-in-hpc-research-exa2pro-dqra-and-hicma-parse-frameworks-more/}{HPCwire}]
    \item KAUST Leverages Mixed Precision for Geospatial Data\hfill[\href{https://www.hpcwire.com/2021/07/28/kaust-leverages-mixed-precision-for-geospatial-data/}{HPCwire}]
    \item Mixing Precision for Model Acceleration\hfill[\href{https://techxplore.com/news/2021-07-precision.html}{Tech Xplore}]
    \item Mixing It Up: Saudi Researchers Accelerate Environmental Models with Mixed Precision\hfill[\href{https://blogs.nvidia.com/blog/2020/04/14/kaust-environmental-models-mixed-precision/}{Nvidia}]
    \item\begin{CJK}{UTF8}{min}「富岳」を用いた3つの研究成果がゴードン・ベル賞ファイナリストに選出されました\end{CJK}\hfill[\href{https://www.r-ccs.riken.jp/outreach/topics/20220826-1/}{RIKEN News}]
  \end{itemize}   
  
\end{rSection}

% footer
\bigskip
\begin{center}
  \begin{footnotesize}
    Last updated: \today
  \end{footnotesize}
\end{center}



%	EXAMPLE SECTION
%----------------------------------------------------------------------------------------

%\begin{rSection}{ACTIVITIES \& AWARDS}
%\item Runners up in B.G.Shirke Vidyarthi C
%\item Won First Prize in Model Making Competition Organized by
%\end{rSection}

\newpage

%----------------------------------------------------------------------------------------
% Extra Curricular
%----------------------------------------------------------------------------------------
%\begin{rSection}{Extra-Cirrucular} \itemsep -3pt
%\item Co-Organized “ Nirmitee 2017” - a National Symposium of Civil Department of MIT, Pune

%\bibliographystyle{unsrt}
%\bibliography{ref}

\end{document}
